\documentclass[11pt,a4paper]{article}
\usepackage[utf8]{inputenc}
\usepackage[T1]{fontenc}
\usepackage[spanish]{babel}
\usepackage[a4paper,margin=2.4cm]{geometry}
\usepackage{amsmath}
\usepackage{booktabs}
\usepackage{array}
\usepackage{xcolor}
\usepackage{listings}
\usepackage{hyperref}
\usepackage{tikz}
\usetikzlibrary{positioning,arrows.meta,shapes.geometric}

\lstset{
  basicstyle=\ttfamily\small,
  breaklines=true,
  frame=single,
  columns=fullflexible,
  showstringspaces=false
}

\hypersetup{colorlinks=true,linkcolor=blue,urlcolor=blue}

\title{\textbf{Proyecto final de Base de Datos -- Parte 2}\\[0.3em]
\Large Arquitectura de datos para Ciencia de Datos e IA:\\[0.2em]
Predicción y detección de anomalías de demanda espaciotemporal de viajes (NYC TLC HVFHV 2026)}
\author{%
  Sebastian Rodrigo García Villacorta\\
  Dante Guillermo Barreto Daza\\
  Clinton Michael Chulluncuy Reynoso\\
  Ronald Ramiro Hilario Orihuela}
\date{Curso: Base de Datos (MCD8008) -- Docente: Luis Alberto Chicoma Chevez}

\begin{document}

\maketitle

\begin{abstract}
En la primera entrega se implementó un modelo relacional en PostgreSQL sobre los registros
High Volume For-Hire Vehicle (HVFHV) de Nueva York (enero--mayo de 2026), con 11 tablas,
particionado mensual, control de calidad y feature engineering orientado a la predicción de
demanda. En esta segunda etapa se amplía el caso hacia una arquitectura especializada
complementaria: se selecciona \textbf{TimescaleDB} (series temporales) junto con
\textbf{PostGIS} (espacial) para responder en tiempo casi real a la pregunta
\emph{¿cuántos viajes se solicitarán en la próxima hora por zona y plataforma, y qué zonas
presentan demanda anómala?}. Se justifica la elección por el patrón de acceso temporal y
espacial de los datos, se diseña el modelo de datos especializado y se presentan resultados
reales de la implementación, incluyendo una reducción de almacenamiento de $162$\,MB a
$14$\,MB mediante compresión nativa (factor $\approx 11{,}6\times$).
\end{abstract}

\tableofcontents
\newpage

\section{Arquitectura de datos para Ciencia de Datos e IA}

El caso desarrollado en la Parte 1 modela los viajes de movilidad bajo demanda de Uber y Lyft
publicados por la \emph{New York City Taxi and Limousine Commission}. El proceso observable
comprende solicitud, asignación, llegada, recogida, desplazamiento, finalización y componentes
económicos de cada viaje.

Las características del problema --- volumen elevado, generación periódica de datos, análisis
predominantemente por rangos de tiempo y por zona geográfica --- hacen que el modelo
relacional sea insuficiente para consultas analíticas de baja latencia sobre series históricas
largas. Por ello se identifica un escenario de Ciencia de Datos que puede beneficiarse de una
tecnología de almacenamiento y procesamiento especializada, \textbf{sin reemplazar} a
PostgreSQL, que sigue siendo la fuente de verdad transaccional.

\section{Definición del caso de IA / Ciencia de Datos}

\subsection{Enunciado}
\textbf{Caso:} predicción de demanda espaciotemporal de viajes y detección de anomalías.

\emph{¿Cuántos viajes se van a solicitar en la próxima hora en cada zona y plataforma, y qué
zonas presentan una demanda inusualmente alta o baja respecto de su comportamiento
habitual?}

\subsection{Unidad predictiva y variable objetivo}
La unidad predictiva, heredada de la Parte 1, es la combinación
\[
(\text{hora},\ \text{zona\_origen\_id},\ \text{plataforma\_codigo}),
\]
y la variable objetivo es $\text{viajes\_objetivo}$, es decir, el número de recogidas observadas
en esa hora, zona y plataforma.

\subsection{Salidas esperadas}
\begin{itemize}
  \item Pronóstico horario de demanda por zona y plataforma (tarea de regresión).
  \item Alerta de anomalía cuando la demanda observada se aleja del comportamiento esperado.
\end{itemize}

\section{Caracterización del problema}

\begin{table}[h]
\centering
\begin{tabular}{@{}p{4.6cm}p{8.6cm}@{}}
\toprule
\textbf{Dimensión} & \textbf{Valor / estimación} \\
\midrule
Volumen & 105\,996\,114 filas en el universo (ene--may 2026); muestra reproducible de
$1\,000\,000$ filas/mes ($4\,999\,931$ cargadas). Proyección $\approx 250$\,M filas/año. \\
Frecuencia & Ingesta mensual por lotes; análisis horario (agregado \texttt{demanda\_hora}). \\
Tipos de datos & Temporales (timestamp), numéricos (importes, millas, duración), booleanos y
geoespaciales (longitud/latitud WGS84). \\
Tiempo real & Casi real: inferencia cada hora sobre la serie actualizada. \\
Latencia esperada & Inferencia $< 100$\,ms; consulta de demanda $< 500$\,ms. \\
Patrón lect/escritura & Escritura en lote (ingesta mensual) y lectura analítica intensiva por
rango temporal y por zona. \\
Escalabilidad & Particionado por tiempo + compresión para sostener $\approx 250$\,M filas/año. \\
Consultas predominantes & Agregados temporales (\texttt{time\_bucket}), rezagos, medias
móviles y filtros por zona/proximidad. \\
\bottomrule
\end{tabular}
\caption{Caracterización del problema de predicción de demanda.}
\end{table}

\section{Selección de tecnología}

Se selecciona \textbf{TimescaleDB} (extensión de PostgreSQL para series temporales) como
tecnología principal, complementada con \textbf{PostGIS} para capacidades espaciales. La
selección se deriva directamente de la caracterización anterior:

\begin{itemize}
  \item \textbf{Qué problema resuelve:} consultas de rango temporal y agregaciones sobre
        millones de observaciones con baja latencia, que en un escaneo secuencial del modelo
        relacional son costosas.
  \item \textbf{Por qué fue seleccionada:} el dato es una serie temporal; TimescaleDB aporta
        hypertables (particionado automático por tiempo), compresión nativa y
        \emph{continuous aggregates}, y al ser una extensión se integra con el modelo
        relacional existente sin migración.
  \item \textbf{Ventajas:} compresión demostrada ($162$\,MB $\to$ $14$\,MB), agregados
        mantenidos automáticamente y consultas espaciales vía PostGIS/GiST.
  \item \textbf{Limitaciones:} la compresión de segmentos deshabilita algunas escrituras
        directas sobre chunks comprimidos (requiere descompresión/recompresión) y el
        paralelismo de escritura es menor que en soluciones distribuidas puras.
  \item \textbf{Cómo se complementa:} PostgreSQL conserva el núcleo transaccional (viajes,
        catálogos, auditoría, integridad referencial); TimescaleDB/PostGIS cubren el patrón de
        acceso analítico temporal--espacial.
\end{itemize}

\section{Modelo de datos NoSQL / especializado}

El modelo especializado se describe en detalle en el entregable \texttt{modelo\_nosql.pdf}.
En síntesis:

\begin{enumerate}
  \item Se añade una columna \texttt{geo GEOGRAPHY(POINT,4326)} a \texttt{core.zona\_taxi},
        derivada de longitud/latitud, con índice espacial GiST.
  \item Se crea el hypertable \texttt{analytics.demanda\_serie} particionado por \texttt{hora}.
  \item Se crea un \emph{continuous aggregate} diario \texttt{analytics.demanda\_diaria}.
  \item Se configura compresión por segmentos (\texttt{zona\_origen\_id},
        \texttt{plataforma\_codigo}) ordenados por \texttt{hora}.
\end{enumerate}

\section{Arquitectura propuesta}

El diagrama completo se presenta en el entregable \texttt{arquitectura.pdf}. Los componentes
principales son:

\begin{center}
\begin{tikzpicture}[
  node distance=0.6cm and 1.2cm,
  box/.style={draw, rounded corners=2pt, align=center, minimum height=0.7cm, minimum width=2.6cm, font=\small},
  arr/.style={-{Stealth[length=2mm]}}
]
\node[box] (app) {APP Web / Mobile};
\node[box, below=of app] (api) {API};
\node[box, below left=of api] (pg) {PostgreSQL\\ (core transaccional)};
\node[box, below right=of api] (ts) {TimescaleDB + PostGIS\\ (series + espacial)};
\node[box, below=of ts] (fs) {Feature Store};
\node[box, below=of fs] (ml) {ML Pipeline\\ (entrenamiento)};
\node[box, below=of ml] (reg) {Model Registry\\ (MLflow)};
\node[box, below=of reg] (inf) {API de inferencia\\ (FastAPI)};
\node[box, below=of pg] (datos) {Datos de negocio};

\draw[arr] (app) -- (api);
\draw[arr] (api) -- (pg);
\draw[arr] (api) -- (ts);
\draw[arr] (ts) -- (fs);
\draw[arr] (fs) -- (ml);
\draw[arr] (ml) -- (reg);
\draw[arr] (reg) -- (inf);
\draw[arr] (pg) -- (datos);
\end{tikzpicture}
\end{center}

\section{Comparación entre tecnologías}

\begin{table}[h]
\centering
\begin{tabular}{@{}p{4.2cm}p{4.6cm}p{4.8cm}@{}}
\toprule
\textbf{Criterio} & \textbf{PostgreSQL} & \textbf{TimescaleDB (+PostGIS)} \\
\midrule
Modelo & Relacional & Serie temporal + espacial \\
Particionado & Manual (por mes) & Automático (hypertable) \\
Compresión & No nativa & Nativa por segmentos ($\approx 11{,}6\times$ observado) \\
Agregados históricos & Vistas/consultas costosas & Continuous aggregates \\
Consultas por rango temporal & Costosas sobre millones de filas & Optimizadas por tiempo \\
Espacial & No & PostGIS + índice GiST \\
Transacciones ACID & Excelente & Excelente (es PostgreSQL) \\
Caso principal & Core transaccional, integridad, catálogos & Analítica temporal--espacial en tiempo casi real \\
\bottomrule
\end{tabular}
\caption{Comparación entre PostgreSQL y TimescaleDB/PostGIS.}
\end{table}

\subsection{¿Por qué utilizar ambas tecnologías en lugar de reemplazar una por otra?}
La respuesta se fundamenta en el concepto de \emph{Polyglot Persistence}: distintas tecnologías
coexisten según las características de cada problema. PostgreSQL permanece como la
\emph{fuente de verdad} transaccional (viajes, catálogos, lotes de ingesta, reglas de calidad e
integridad referencial). TimescaleDB y PostGIS optimizan el patrón de acceso analítico
temporal--espacial que el modelo relacional no resuelve eficientemente. Los datos fluyen de
PostgreSQL hacia la capa especializada, sin duplicar la lógica de negocio ni sustituir el
sistema de registro.

\section{Resultados de la implementación}

La implementación se ejecutó sobre PostgreSQL 16 + TimescaleDB 2.31 + PostGIS 3.6 en un
contenedor Docker. Los resultados son reproducibles con los scripts de la carpeta
\texttt{Sql/}.

\subsection{Caracterización}
\begin{itemize}
  \item Viajes cargados: $4\,999\,931$ (Uber $3\,590\,362$ = $71{,}81\%$; Lyft $1\,409\,569$
        = $28{,}19\%$).
  \item Agregado horario relacional \texttt{demanda\_hora}: $1\,272\,029$ filas ($162$\,MB).
  \item Hypertable \texttt{demanda\_serie}: $1\,272\,029$ filas, $22$ chunks, todos
        comprimidos: $14$\,MB (factor $\approx 11{,}6\times$).
\end{itemize}

\subsection{Consultas temporales}
\begin{itemize}
  \item Hora pico global: \texttt{2026-02-07 19:00} con $3\,163$ viajes.
  \item Día de mayor demanda: \texttt{2026-02-07} con $47\,981$ viajes.
\end{itemize}

\subsection{Consultas espaciales}
\begin{itemize}
  \item Zonas con mayor demanda en un radio de $2$\,km de Midtown Center: Times Sq
        ($55\,043$), Midtown Center ($55\,014$), Midtown South ($45\,399$).
  \item Zonas más cercanas a JFK (kNN): Springfield Gardens South ($2\,514$\,m), Springfield
        Gardens North ($3\,012$\,m), Baisley Park ($3\,570$\,m).
  \item Distancia de centroides Midtown--JFK: $20\,324$\,m ($12{,}63$ millas), frente a
        $17{,}56$ millas de media en los $1\,040$ viajes reales con esa ruta.
\end{itemize}

\section{Pasos de replicación}

\begin{enumerate}
  \item Levantar el stack con \texttt{docker compose up -d --build} (carpeta \texttt{deploy/}).
  \item Ejecutar en orden \texttt{01\_schema.sql}, \texttt{02\_tables.sql},
        \texttt{03\_data.sql} y \texttt{06\_zonas\_geo.sql}.
  \item Cargar los CSV de \texttt{viaje} e \texttt{incidencia} con
        \texttt{\textbackslash copy ... WITH (FORMAT csv, HEADER true, NULL '')}.
  \item Reconstruir \texttt{analytics.demanda\_hora}.
  \item Ejecutar \texttt{07\_timescale.sql} (capa espacial + hypertable + continuous
        aggregate + compresión).
  \item Ejecutar \texttt{08\_demo\_queries.sql} para reproducir los resultados.
\end{enumerate}

\end{document}
